% Capitulo 1: Introduccion (Proyecto de Tesis -- Fase 8, docs/rio_search_plan.md S3.11)
\chapter{Introducci\'on}

\section{Problema y motivaci\'on}

Este proyecto de tesis propone determinar qu\'e metodolog\'ia -- modelo, ventana de entrenamiento,
conjunto de \emph{features} y estrategia de horizontes -- proyecta mejor el caudal del r\'io
Uruguay en su cuenca alta, con el punto de predicci\'on fijado en la estaci\'on
\texttt{ana\_74100000} (Decisi\'on 018, \texttt{docs/decisions.md}). El problema es relevante para
la gesti\'on del recurso h\'idrico -- alerta de crecidas, planificaci\'on de infraestructura y
actividades productivas aguas abajo -- y se aborda como un problema de series temporales
multi-horizonte (1 a 7 y 14 d\'ias) sobre datos hidrol\'ogicos y meteorol\'ogicos observados.

\section{Objetivos}

\subsection{Objetivo general}

Construir un banco de pruebas reproducible (\emph{Rio\_Search}) que entrene, compare y opere
metodolog\'ias de proyecci\'on de caudal sobre \texttt{weather.gold.training\_dataset\_v0}, con
trazabilidad completa (dataset, configuraci\'on, c\'odigo, m\'etricas y tiempos) en MLflow.

\subsection{Objetivos espec\'ificos}

\begin{itemize}[noitemsep]
  \item Establecer un piso de comparaci\'on con modelos \emph{naive} (persistencia, climatolog\'ia,
    estacionalidad naive).
  \item Entrenar y evaluar un modelo BiLSTM bajo dos estrategias de horizonte
    (\texttt{multi\_output} y \texttt{per\_horizon}) y comparar su costo computacional.
  \item Dise\~nar el protocolo de evaluaci\'on hidrol\'ogica (NSE, KGE, \emph{skill} vs.
    persistencia) y de selecci\'on de modelo campe\'on sobre validaci\'on, no sobre prueba.
  \item Dejar el dise\~no abierto a incorporar el pron\'ostico meteorol\'ogico (ECMWF/GEFS) y
    modelos adicionales sin tocar el n\'ucleo del dominio.
\end{itemize}

\section{Alcance}

El alcance de este proyecto de tesis cubre exclusivamente la cuenca alta (Decisi\'on 018); el
segundo punto de predicci\'on hist\'oricamente considerado (aguas abajo, asociado a Salto Grande)
queda expl\'icitamente fuera. El proyecto no reemplaza el pipeline de ingesta y consolidaci\'on de
datos (\texttt{docs/roadmap.md}), que sigue su propio curso independiente; \emph{Rio\_Search}
consume lo que ese pipeline produce.
